\section{Structure of programs}
%\label{sec:4:ecrir-progs}

Just as writing correct English requires following precise rules of
spelling and grammar, writing correct programs is subject to a number
of {\em syntax rules}\index{Syntax rule}, fortunately simpler than
those of English. The difference
is that English sentences containing spelling or grammar mistakes
may still be understood by English speakers, whereas
a program containing
{\em syntax errors}\index{Syntax error} cannot
be translated into the language of the com\-pu\-ter, and therefore cannot be run.
This translation is carried out by a program (called an {\em interpreter})\index{Interpreter}
designed to point out our errors but not to correct them. To correct
our errors in a satisfactory way, the interpreter would have to understand
our intentions, a very difficult problem.

The program {\idt PaperHen} below is an example of a
correct \python\ program.

\begin{exemple}[Syntactically correct program]{\hfill
\label{prog:Cocotte} \input python/PaperHen.tex
\noindent $\Box$ }\end{exemple}

\noindent All our \python\ programs for driving the \robot\ have the same form:

\begin{center}
\begin{tabular}{@{}|p{12.5cm}|@{}} \hline
{\em \python\ program:}\\ \hline \hline
\begin{minipage}[t]{12.5cm}
{\medskip
\mbox{\blc\ \ \ {\mcl from}\hspace{0.3em}\blanc\hspace{0.3em}{\idt LibRobot}\hspace{0.3em}\blanc\hspace{0.3em}{\mcl import}\hspace{0.3em}\blanc\hspace{0.3em}{\idt *}}\par
\mbox{\blc\ \ \ \cmt \#\ \small\rm Comment explaining what the program does}\par
\mbox{\blc\ \ \ \small\rm  \ldots\ Instructions of the program (one per line, or separated by `$;$') \ldots}\par
\medskip}
\end{minipage}\\ \hline
\end{tabular}
\end{center}

A \python\ program is stored in a file whose name ends with {\idt .py}.
The name of the program is the name of this file, without the {\idt .py}: the program
{\idt PaperHen} is stored in the file {\idt PaperHen.py}.

This name is chosen by the programmer
as he or she likes, so as to suggest what the program is for.
 The only rule to fol\-low is that the name of the program
must contain only unaccented lowercase or
uppercase letters\footnote{\ Accents do not exist in English
(with very rare exceptions, as in the words `naïve', `protégé' or
`rôle', borrowed from French).}, digits or
underscores `\_'\footnote{\ The underscore `\_' must not be confused with the hyphen `-'.}
and must begin with a letter. For example {\idt  S},
{\idt  TheStegosaurus} or {\idt  The\_Stegosaurus} are correct program
names whereas {\idt  1S},
{\idt  s!}, {\idt  The\ Stegosaurus} or {\idt  The-Stégosaurus} do not follow
the rule for writing program names.

A \python\ program begins with the line `{\mcl from}\ {\idt LibRobot}\ {\mcl import}\ {\idt *}'.
This line
imports from the
{\em library}\index{Library} {\idt LibRobot} all (that is the meaning of the star `{\idt *}')
the commands of the
\robot, without which the computer would not know them. It also opens the window
in which the \robot\ draws.

There must be a blank
between the word {\mcl from} and the name {\idt LibRobot}, as well as between {\idt LibRobot} and the
word {\mcl import}. In a book there are blanks everywhere and even, if you look closely,
more blank than anything else. Putting a blank in a book to say that you must type
a blank is therefore silly, because most readers are not interested in the
blanks in books but simply in what is written there.
We shall therefore use the symbol `\blanc' to write our blanks in black and \hbox{white!}
In other words we write \blanc\ to say that you must press the long bar
at the bottom of the com\-pu\-ter keyboard, called the
{\em space bar}\index{Space bar},
to obtain a blank `\ ' in the program.
We do not do so in the programs given as examples, which are
shown as they appear on the screen. Likewise, at the end of a
line you must press the return key
\framebox[11.mm][l]{$\hookleftarrow$~~~}\hspace{0.3ex} or \hspace{0.3ex}\framebox[11.mm][l]{~~\tiny\texttt{return}}\hspace{0.3ex} to go to the beginning of the
next line. We do not write it in books because all readers know
that they must go to the beginning of the next line when they reach the end of a line.

After the comment that explains what the program does
come the instructions
of the program. They are written in the order in which they must be executed by
the com\-pu\-ter, starting at the beginning of the line (with no blank in front: we shall see
in the next chapter that blanks at the beginning of a line have a precise meaning in \python).
You can write one instruction per line or else several instructions on the same line,
separating them by semicolons, with no semicolon after the last one on the line.
The instructions of the program
can be commands of the \hbox{\robot:} {\idt fd}, {\idt cl}, {\idt cr}, {\idt qtr},
{\idt qtl}, {\idt etr}, {\idt etl}, {\idt pu}, {\idt pd}, {\idt dc}, {\idt dp},
{\idt dg} or {\idt st}, each followed by {\idt ()}. In general, programs for driving the \robot\
end with the command {\idt st()}, which waits until we
press a key before the drawing
disappears at the end of the execution of the program.

Texts that begin with the sign {\cmt \#} and continue to the end of
the line are {\em comments}, ignored by the computer.

The pro\-gram
{\idt WrongPaperHen} below contains
several errors that are very common among beginners.

\begin{exemple}[Syntactically incorrect program]{\hfill

\label{prog:Cocotte-errone} \input python/WrongPaperHen.tex
\noindent The errors in this program are the \hbox{following:}
\begin{list}{$\scriptscriptstyle\bullet$}{
\renewcommand{\labelwidth}{3mm}
\renewcommand{\labelsep}{3mm}
\renewcommand{\leftmargin}{10mm}
\renewcommand{\rightmargin}{0mm}
}
\item[1.] the line `{\mcl from}\ {\idt LibRobot}\ {\mcl import}\ {\idt *}' has been forgotten;
\item[2.] the text that explains what the program does has been written between triple
quotes {\cha \char34\char34\char34} (it is a {\em documentation string}, which can
be used as a comment spanning several lines) but the closing triple quotes
{\cha \char34\char34\char34} have been forgotten;
\item[3.] the instruction {\idt rtl()} is not a command of the \robotErr;
\item[4.] the opening parenthesis of the last command {\idt st(} is not closed.
\end{list}
\noindent $\Box$ }\end{exemple}


A complicated example of a syntactically correct \python\ program
is given \hbox{below:}
\begin{exemple}[The stegosaurus]{\rm
The \python\ program for driving the \robot\ to
draw the ste\-go\-saur\-us shown in fig\-ure
\ref{fig:stegosaure} (page \pageref{fig:stegosaure}) is the \hbox{following:}
\label{prog:Stegosaure} \input python/Stegosaurus.tex
\noindent $\Box$ }\end{exemple}

\section{Comments}
%\label{sec:4:comment}

A \python\ program may contain {\em comments}\index{Comment}, which are
any texts written after the sign {\cmt \#\ \ldots}, up to the end of the
line. A comment may take up a whole line or follow the instructions
of a line, as in {\idt cl()\ \ }{\cmt \#\ ERROR:\ cr}.
A text of several lines may also be placed between triple quotes
{\cha \char34\char34\char34\ \ldots\ \char34\char34\char34}: it is a
{\em documentation string}\index{Documentation string}, which the computer
reads but which gives no command to the \robot.
The com\-pu\-ter takes no account of these comments, which are only
meant to explain the use and the working principle of the
program to the readers of this program. If all the
comments of a program are removed, it becomes unreadable for a human but
nothing is changed for the com\-pu\-ter.

\begin{exemple}[The \underline{unreadable} stegosaurus]{\rm
By removing all the com\-ments of the program {\idt Stegosaurus}
on page \pageref{prog:Stegosaure}, as well as the useless blanks and line
breaks, we obtain a \python\ program that is
totally unreadable but which draws perfectly well the
ste\-go\-saur\-us shown in fig\-ure
\ref{fig:stegosaure} (page \hbox{\pageref{fig:stegosaure}):}
{\small{
\parindent \parindentprog
\parskip 0pt
\nopagebreak\medskip

\carrenoir
\vspace{\neglskip}\mbox{\blc {\mcl from}{\idt \ LibRobot\ }{\mcl import}{\idt \ *}}\par
\nopagebreak
\vspace{\neglskip}\mbox{\blc {\idt pu();\ qtr();\ qtr();\ fd();\ fd();\ fd();\ fd();\ fd();\ fd();\ fd()}}\par
\vspace{\neglskip}\mbox{\blc {\idt qtl();\ fd();\ fd();\ fd();\ fd();\ fd();\ fd();\ fd();\ fd();\ fd()}}\par
\vspace{\neglskip}\mbox{\blc {\idt fd();\ fd();\ fd();\ fd();\ qtl();\ pd();\ fd();\ qtr();\ cl();\ qtl()}}\par
\vspace{\neglskip}\mbox{\blc {\idt cl();\ qtr();\ qtr();\ cl();\ qtr();\ qtr();\ cl();\ qtl();\ cl()}}\par
\vspace{\neglskip}\mbox{\blc {\idt qtr();\ cr();\ cl();\ fd();\ cr();\ fd();\ fd();\ fd();\ qtr();\ fd()}}\par
\vspace{\neglskip}\mbox{\blc {\idt qtr();\ fd();\ fd();\ fd();\ fd();\ qtr();\ cr();\ fd();\ qtl();\ cr()}}\par
\vspace{\neglskip}\mbox{\blc {\idt fd();\ fd();\ fd();\ qtr();\ fd();\ qtr();\ fd();\ fd();\ qtr();\ pu()}}\par
\vspace{\neglskip}\mbox{\blc {\idt fd();\ pd();\ cr();\ fd();\ fd();\ qtr();\ fd();\ qtr();\ fd();\ qtr()}}\par
\vspace{\neglskip}\mbox{\blc {\idt pu();\ fd();\ pd();\ cr();\ qtl();\ fd();\ qtr();\ fd();\ qtr();\ fd()}}\par
\vspace{\neglskip}\mbox{\blc {\idt fd();\ qtr();\ fd();\ fd();\ qtr();\ cr();\ qtl();\ fd();\ fd();\ cr()}}\par
\vspace{\neglskip}\mbox{\blc {\idt fd();\ qtr();\ fd();\ fd();\ qtr();\ fd();\ fd();\ qtr();\ fd();\ cr()}}\par
\vspace{\neglskip}\mbox{\blc {\idt qtl();\ fd();\ fd();\ qtr();\ fd();\ fd();\ qtr();\ fd();\ fd()}}\par
\vspace{\neglskip}\mbox{\blc {\idt qtr();\ fd();\ fd();\ qtr();\ fd();\ fd();\ qtl();\ cr();\ fd()}}\par
\vspace{\neglskip}\mbox{\blc {\idt qtr();\ fd();\ fd();\ qtr();\ fd();\ fd();\ qtr();\ fd();\ cr();\ fd()}}\par
\vspace{\neglskip}\mbox{\blc {\idt fd();\ qtl();\ cr();\ qtr();\ fd();\ fd();\ qtr();\ fd();\ fd()}}\par
\vspace{\neglskip}\mbox{\blc {\idt qtr();\ fd();\ qtr();\ fd();\ qtl();\ cr();\ cl();\ qtr();\ qtr()}}\par
\vspace{\neglskip}\mbox{\blc {\idt fd();\ fd();\ qtr();\ fd();\ qtr();\ fd();\ qtr();\ cl();\ cr();\ cl()}}\par
\vspace{\neglskip}\mbox{\blc {\idt qtr();\ qtr();\ fd();\ fd();\ fd();\ qtr();\ fd();\ fd();\ qtr()}}\par
\vspace{\neglskip}\mbox{\blc {\idt fd();\ cr();\ cl();\ fd();\ cr();\ qtr();\ fd();\ fd();\ fd();\ qtr()}}\par
\vspace{\neglskip}\mbox{\blc {\idt fd();\ qtr();\ fd();\ fd();\ cr();\ fd();\ fd();\ pu();\ qtl();\ fd()}}\par
\vspace{\neglskip}\mbox{\blc {\idt dc();\ qtr();\ qtr();\ fd();\ qtl();\ pd();\ cl();\ cr();\ fd();\ cl()}}\par
\vspace{\neglskip}\mbox{\blc {\idt qtl();\ fd();\ fd();\ qtr();\ qtr();\ fd();\ fd();\ qtl();\ cl()}}\par
\vspace{\neglskip}\mbox{\blc {\idt fd();\ fd();\ fd();\ fd();\ fd();\ fd();\ fd();\ cl();\ qtr();\ fd()}}\par
\vspace{\neglskip}\mbox{\blc {\idt qtl();\ fd();\ fd();\ fd();\ qtr();\ qtr();\ fd();\ fd();\ fd();\ fd()}}\par
\vspace{\neglskip}\mbox{\blc {\idt fd();\ qtl();\ fd();\ fd();\ qtl();\ fd();\ fd();\ fd();\ fd();\ fd()}}\par
\vspace{\neglskip}\mbox{\blc {\idt fd();\ qtr();\ qtr();\ fd();\ fd();\ fd();\ fd();\ fd();\ qtl();\ fd()}}\par
\vspace{\neglskip}\mbox{\blc {\idt qtl();\ fd();\ fd();\ qtl();\ fd();\ qtr();\ qtr();\ fd();\ cl()}}\par
\vspace{\neglskip}\mbox{\blc {\idt qtr();\ fd();\ fd();\ fd();\ qtl();\ fd();\ fd();\ fd();\ qtr()}}\par
\vspace{\neglskip}\mbox{\blc {\idt qtr();\ fd();\ fd();\ fd();\ fd();\ fd();\ fd();\ fd();\ qtl();\ fd()}}\par
\vspace{\neglskip}\mbox{\blc {\idt fd();\ qtl();\ fd();\ fd();\ fd();\ fd();\ fd();\ fd();\ fd();\ fd()}}\par
\vspace{\neglskip}\mbox{\blc {\idt qtr();\ qtr();\ fd();\ fd();\ fd();\ fd();\ fd();\ fd();\ fd();\ qtl()}}\par
\vspace{\neglskip}\mbox{\blc {\idt fd();\ qtl();\ fd();\ fd();\ fd();\ qtl();\ fd();\ qtr();\ qtr()}}\par
\vspace{\neglskip}\mbox{\blc {\idt fd();\ fd();\ qtr();\ cl();\ fd();\ cr();\ cl();\ fd();\ qtr();\ cl()}}\par
\nopagebreak
\vspace{\neglskip}\mbox{\blc {\idt qtr();\ st()}}\par

\nopagebreak\medskip }\nopagebreak
}
This example also shows that the layout of the instructions of the
program is important. Spacing between instructions and
proper alignment make programs much easier
to read.
\noindent $\Box$ }\end{exemple}



\section{Interpretation of programs}
%\label{sec:4:compil-progs}
The {\em interpretation}\index{Interpretation} of a \python\ program consists in translating it
into the language of the com\-pu\-ter,
called {\em machine language}\index{Machine language}, and having it executed by the computer.
It is done by a program
called the \index{Interpreter}\python\ {\em interpreter}\@\footnote{\ The command
{\tt ./robot.sh PaperHen.py}, typed in a
terminal window opened in the folder {\tt LeRobot\_Python\_EN}, calls the interpreter
{\tt python3} to run the program {\tt PaperHen.py}. Dragging the icon of
{\tt PaperHen.py} onto that of the application {\tt LeRobotPythonEN.app} does the same thing.}.
The interpreter works in two stages:
\begin{list}{$\scriptscriptstyle\bullet$}{
\renewcommand{\rightmargin}{0mm}
}
\item it first reads the {\em whole} program to check that it follows the syntax rules
of \python\ and to translate it; if it finds a syntax error, it reports it
and stops without running anything;
\item then, if the program is syntactically correct, it has the computer execute its
instructions, one after the other.
\end{list}
Some errors can be discovered only during this second stage, when
the com\-pu\-ter reaches the faulty instruction, for example the call of a command that does not
exist. Execution then stops, but the \robot\ has already made all the drawing requested
by the preceding instructions.
Obviously the \python\ program and its translation drive
the \robot\ in the same way. This way of proceeding
saves us from having to learn machine language. It is very difficult
to understand this machine language because all programs are written
only with 0s and \hbox{1s!}

\section{Syntax errors}
%\label{sec:4:err-compil}

To be able to run a \python\ program, it must not contain
any syntax errors\index{Programming error}\index{Syntax error} violating the rules
for writing programs. In case of an error the interpreter explains
briefly, in English, the cause of the error and points out the place in the program
where the er\-ror was found (the line number, for example {\tt line 15}, then the line
itself and, underneath, one or more
carets {\idt \char94} under the faulty text). The message begins with
the name of the error, for example {\tt SyntaxError}. You must then
correct the error and run the program again.

An error discovered during execution is reported in the same way, but the
message begins with {\tt Traceback (most recent call last):}
and ends with the name of the error, for example {\tt NameError}.

\begin{exemple}[Syntactically incorrect program]{\rm
\hfill The pro\-gram {\idt WrongPaperHen} on page
\pageref{prog:Cocotte-errone}
contains several of the most frequent errors for
such a simple program.
When we run this program {\idt WrongPaperHen}, the interpreter starts by reporting
`{\tt SyntaxError: unterminated triple-quoted string literal (detected at line 15)}',
pointing at
the {\cha \char34\char34\char34} of the
first line: the closing triple quotes were forgotten at the end of the
text that explains what the program does, so that the whole rest of the
program, up to the last line 15, is part of this text!
Having added the forgotten {\cha \char34\char34\char34} at the end of the first line,
a new attempt reports
`{\tt SyntaxError: \textquotesingle{}(\textquotesingle{} was never closed}'
pointing at the parenthesis of {\idt st(} on line 15. We close it by writing {\idt st()}.
The program is now syntactically correct and its execution can begin.
It stops at once, on the first instruction, with the message
`{\tt NameError: name \textquotesingle{}etr\textquotesingle{} is not defined. Did you mean: \textquotesingle{}str\textquotesingle{}?}'
(the suggestion {\idt str}, a name known to \python, is of no help here).
Since the line `{\mcl from}\ {\idt LibRobot}\ {\mcl import}\ {\idt *}' was forgotten,
the interpreter does not know the meaning of the commands of the \robotErr, and the window
of the \robot\ has not even opened.
After adding this line at the beginning of the program, we can run it again.
The window opens, the \robot\ draws the first line of the back of the paper hen,
then everything stops with the message
`{\tt NameError: name \textquotesingle{}rtl\textquotesingle{} is not defined. Did you mean: \textquotesingle{}etl\textquotesingle{}?}',
which points at the
line where {\idt rtl()} is. It is a typing
mistake, which we correct by replacing {\idt  rtl} by {\idt  etl}.
We thus laboriously obtain the program {\idt PaperHen} on page
\pageref{prog:Cocotte}, which is correct (apart from the comment, written with {\cmt \#} in {\idt PaperHen}).
The interpreter stops at the first error it finds: you must therefore
run the program again after each correction. Note that the two syntax
errors were reported before any execution, whereas the two unknown names
were discovered only during execution, when the com\-pu\-ter reached
the faulty instruction.
\noindent $\Box$ }\end{exemple}


\section{Execution of programs}
%\label{sec:4:exec-progs}
The {\em execution}\index{Execution} of a program consists in having
the com\-pu\-ter carry out the actions defined by the instructions of the program. Since in our case the instructions are
commands of the \robot, executing the program makes it possible to drive the
\robot.

{\newcommand{\touche}[1]{{\begin{picture}(12,12)(0, 2)\put(0,0){\framebox(12,12)[lb]{}}\put(1.5,2){\makebox(0,0)[lb]{\tiny #1}}\end{picture}\hspace{0.3ex}}}
When the \robot\ has finished its drawing, the command {\idt st()} displays at the bottom of
the window the keys that let you \hbox{continue:}
\begin{list}{$\scriptscriptstyle\bullet$}{
\renewcommand{\rightmargin}{0mm}
}
\item decrease (by pressing the minus key \touche{$-$} once or several times)
or increase (by pressing the plus key \touche{$+$} once or several times)
the speed of the \robot\ (which can vary from 0 to 10);
\item have the \robot\ redraw the whole drawing, at the chosen speed, by pressing
another key, for example the return key \framebox[11.mm][l]{$\hookleftarrow$~~~}\hspace{0.3ex} or \hspace{0.3ex}\framebox[11.mm][l]{~~\tiny\texttt{return}}\hspace{0.3ex};
when the \robot\ draws too fast for you to watch
the drawing un\-fold, just decrease its speed and then have it
redraw the drawing;
\item stop the execution of the program for good by pressing the
escape key\index{Escape key} \touche{Esc} (or one of the keys
\touche{\S} or \touche{\#}), which gives the screen of figure
\ref{fig:ecran-arret} (page \pageref{fig:ecran-arret}).
\end{list}

\index{Execution (interruption)}\index{Interrupting execution}During
the execution of a program driving the \robot, we
can also stop the drawing of the \robot\ at any moment by pressing the
escape key \touche{Esc} (or \touche{\S} or \touche{\#}).
The drawing then stays displayed as it was at that instant, until we press
a key. For example if we interrupt the execution
of the program {\idt Stegosaurus} on page \pageref{prog:Stegosaure}
when the \robot\ finishes drawing the left front leg,
we obtain a screen similar to the one in figure
\ref{fig:ecran-stop} on page \pageref{fig:ecran-stop}.
}

\begin{figure}[p] %[htbp]
 \includevspacePostScript[7.6cm]{ecran-stop}{0.95}{0.95}
 \caption{Com\-pu\-ter screen after a stop}\label{fig:ecran-stop}
\end{figure}

\begin{figure}[p] %[htbp]
 \includevspacePostScript[7.6cm]{ecran-arret}{0.95}{0.95}
 \caption{Com\-pu\-ter screen before the final stop}
 \label{fig:ecran-arret}
\end{figure}

\section{Logical errors}
%\label{sec:4:err-exec}

\index{Logical error}\index{Programming error}Having eliminated all the syntax errors, that is,
the programming errors that come from violating the grammar rules
of the \python\ language, our troubles are not yet over before we can
admire the \hbox{drawing!} There remain the
{\em logical errors},
which we shall discover when running the program.

The first possible error is that the sequence of commands that our
program asks the com\-pu\-ter to give to the \robot\ leads it
outside the frame. This is the case of the following program:
\input python/LeaveTheFrame.tex
\noindent This program ({\idt LeaveTheFrame}) gives the sad result
below, accompanied in the terminal window by the message
`{\tt Error line 8 of LeaveTheFrame.py: the robot would leave the frame (fd)}':
%indiqué à la figure \ref{fig:ecran-sortir} (page \pageref{fig:ecran-sortir}).
\begin{figure}[htbp]
 \includevspacePostScript[7.6cm]{ecran-sortir}{0.95}{0.95}
 \caption{Com\-pu\-ter screen when the robot leaves the frame}
 %\label{fig:ecran-sortir}
\end{figure}

The second possible error is that we made a mistake in
the program and that the drawing obtained does not look at all like what
we had imagined. For example, the program {\idt WrongStegosaurus}
below is a logically wrong but syntactically
correct version of the program {\idt Stegosaurus} on page \pageref{prog:Stegosaure}.
\input python/WrongStegosaurus.tex
The author of this program {\idt WrongStegosaurus} twice mixed up his left
and his right. The
result of running this program, as
shown
below
(figure \ref{fig:StegosaureErrone}), %page \pageref{fig:StegosaureErrone}
would sur\-prise more than one pa\-lae\-on\-tol\-o\-gist!
\begin{figure}[htbp]
 \includehvspacePostScript[105mm,6.8cm]{StegosaureErrone}{0.95}{0.95}
 \caption{A rather botched stegosaurus}
 \label{fig:StegosaureErrone}
\end{figure}

A third error, typical of \python\ beginners, is to forget the parentheses
{\idt ()} after a command of the \robot, for example `{\idt fd;\ qtr()}' instead of
`{\idt fd();\ qtr()}'. The program is syntactically correct but the \robot\ does not move forward,
because {\idt fd} without {\idt ()} names the command without asking for it to be executed.
The \python\ interpreter reports nothing at all: no error, not even
a {\em warning}\index{Warning}.
Only the drawing, which is missing a line, reveals the error. So you must always
check that every command is indeed followed by its parentheses.

\section{How to locate logical errors in the program}
%\label{sec:4:trouv-err-exec}

To find the logical errors in the program you have to imagine
that you are executing the program in place of the \robotErr. When the program
is very long this is not easy.
The best way is therefore to follow the execution of the program in {\em step-by-step} mode,
\index{Step-by-step execution}
which lets you follow the commands transmitted by the com\-pu\-ter to the \robotErr\ so as to locate the error precisely. When the error is
located, you must stop the execution
of the program in order to correct it in the program. Once
this is done, you can run the program again for a new
attempt.
{\newcommand{\touche}[1]{{\begin{picture}(12,12)(0, 2)\put(0,0){\framebox(12,12)[lb]{}}\put(1.5,2){\makebox(0,0)[lb]{\tiny #1}}\end{picture}\hspace{0.3ex}}}
\begin{list}{$\scriptscriptstyle\bullet$}{
\renewcommand{\rightmargin}{0mm}
}
\item To switch to step-by-step mode we \hbox{can:}
\end{list}
\begin{list}{-}{
\renewcommand{\labelwidth}{3mm}
\renewcommand{\labelsep}{3mm}
\renewcommand{\leftmargin}{15mm}
\renewcommand{\rightmargin}{0mm}
}
  \item at the end of the drawing, reduce the speed to zero by pressing the key
  \touche{$-$} several times, then have the drawing redrawn by pressing the return key
  \framebox[11.mm][l]{$\hookleftarrow$~~~}\hspace{0.3ex} or \hspace{0.3ex}\framebox[11.mm][l]{~~\tiny\texttt{return}};
  \item or insert an instruction {\idt ss()} in the \python\ program, a
  little before the place where the error is suspected. This instruction {\idt ss()}
  (\underline{s}ingle \underline{s}tep) switches
  the com\-pu\-ter to step-by-step execution mode.
\end{list}
\begin{list}{$\scriptscriptstyle\bullet$}{
\renewcommand{\rightmargin}{0mm}
}
\item In step-by-step mode, the com\-pu\-ter waits for us to give it
the order to execute each command of the \robot\ by pressing a key, for example the
return key \framebox[11.mm][l]{$\hookleftarrow$~~~}\hspace{0.3ex} or \hspace{0.3ex}\framebox[11.mm][l]{~~\tiny\texttt{return}}. We can thus follow
the execution of the program instruction after
instruction, which in general makes it fairly easy to find the error.
\enlargethispage{2\baselineskip}% touches return ajoutees page 50 : garder la fin de la liste sur la page
\item When the error has been precisely located in the program, you
 can stop the execution of the program by pressing
 the es\-cape key \touche{Esc} (or \touche{\S} or \touche{\#}).
\item[]  It is also possible to program the return to a normal
   \underline{ex}e\-cu\-tion speed by inserting an instruction {\idt ex()}
  in the program.
\end{list}
}
It almost always takes several attempts before the program is
lo\-gi\-cal\-ly correct.
When it is, one is often tempted to improve it. The hardest thing
is then to know when to \hbox{stop!}

\begin{table}[htbp]
\begin{center}
\begin{tabular}{@{}|l|p{11cm}|@{}} \hline
{\idt st()} & \underline{st}ops the execution of the program at the end of the drawing.
             The speed can be decreased with
 \begin{picture}(12,12)(0, 2)\put(0,0){\framebox(12,12)[lb]{}}\put(1.5,2){\makebox(0,0)[lb]{$\scriptscriptstyle -$}}\end{picture}\hspace{0.3ex} (minimum 0) or increased with
 \begin{picture}(12,12)(0, 2)\put(0,0){\framebox(12,12)[lb]{}}\put(1.5,2){\makebox(0,0)[lb]{$\scriptscriptstyle +$}}\end{picture}\hspace{0.3ex} (maximum 10).
			 Any other key has the drawing redrawn at the indicated speed.
			 Execution stops for good when you press
			 \begin{picture}(12,12)(0, 2)\put(0,0){\framebox(12,12)[lb]{}}\put(1.5,2){\makebox(0,0)[lb]{\tiny esc}}\end{picture}\hspace{0.3ex} or
  \begin{picture}(12,12)(0, 2)\put(0,0){\framebox(12,12)[lb]{}}\put(1.5,2){\makebox(0,0)[lb]{\tiny \S}}\end{picture}\hspace{0.3ex} or
  \begin{picture}(12,12)(0, 2)\put(0,0){\framebox(12,12)[lb]{}}\put(1.5,2){\makebox(0,0)[lb]{\tiny \#}}\end{picture}\hspace{0.3ex}.\\ \hline
{\idt ss()} & executes the rest of the program in step-by-step mode
             (\underline{s}ingle \underline{s}tep, speed 0): each command of the \robot\
			is executed when you press a key.\\ \hline
{\idt ex()} & executes the rest of the program in normal \underline{ex}ecution
             mode (speed 10). \\ \hline
\end{tabular}
\end{center}
\label{tab:cmd-vit1}
\caption{Commands controlling the driving mode of the \protect\robot}
\end{table}

\needspace{6\baselineskip}\noindent \rule{137mm}{0.1mm}\par\nopagebreak %---------------------------------------------
\begin{exercices}{\ \\

\vspace{-\baselineskip}
\noindent \begin{minipage}[t]{88mm}
\tiret Write, debug and run \python\ programs drawing the
\robot\ as it is stylized opposite, with its
pen down and with its pen up.
\end{minipage}
\hfill
\begin{minipage}[t]{44mm}
\vspace{2.0cm}\includePostScript{Robot+BC+LC}{0.95}{0.95}
\hspace{-4.4cm}  % 4.4cm de large
\end{minipage}

\noindent \begin{minipage}[t]{103mm}
\tiret  Write, debug and run a \python\ program that
drives the \robot\ to make it draw the dog shown
opposite.
\end{minipage}
\hfill
\begin{minipage}[t]{29mm}
\vspace{2.3cm}\includePostScript{Chien}{0.95}{0.95}
\hspace{-2.9cm}  % 2.9cm de large
\end{minipage}

\noindent \begin{minipage}[t]{103mm}
\tiret  Program the drawings on
pa\-ge \pageref{fig:quadrilateres}.
\end{minipage}

\vspace{1.5mm}

\tiret  Program the drawing of your initials in capital letters
in the style of the letter {\tt H} in relief shown in
figure \ref{fig:h-relief} (page \pageref{fig:h-relief}).

\noindent $\Box$}\end{exercices}

\begin{corrige}{We give one program ({\idt Dog}), the others being in the folder {\idt book/\allowbreak 04\_writing\_running}.
\input python/Dog.tex
\noindent $\Box$}\end{corrige}

\vfill

%%%%%%%%%%%%%%%%%%%%%%%%%%%%%%%%%%%%%%%%%%%%%%%%%%%%%%%%%%%%%%%%%%%%%%%%%%%
\noindent \framebox[\textwidth]{
\begin{minipage}{\myparboxwidth}
\setlength{\parindent}{2em}
\small \sf
\vspace{\abovedisplayskip}
\section*{Addenda}
The cycle of writing, syntactic debugging, translation, logical
debugging and execution of programs has become considerably
shorter in recent years thanks to the fantastic progress
in computer architecture, which should be understood very
schematically.
A micro\-com\-pu\-ter \hbox{comprises:}
\begin{list}{$\scriptscriptstyle\bullet$}{
}
\item an arithmetic and logic unit to do calculations and
handle data,
\item a memory to keep data
and programs temporarily (its content is lost when the power is switched off),
\end{list}

\vspace{\abovedisplayskip}
\end{minipage}
}
%%%%%%%%%%%%%%%%%%%%%%%%%%%%%%%%%%%%%%%%%%%%%%%%%%%%%%%%%%%%%%%%%%%%%%%%%%%

%%%%%%%%%%%%%%%%%%%%%%%%%%%%%%%%%%%%%%%%%%%%%%%%%%%%%%%%%%%%%%%%%%%%%%%%%%%
\noindent \framebox[\textwidth]{
\begin{minipage}{\myparboxwidth}
\setlength{\parindent}{2em}
\small \sf
\vspace{\abovedisplayskip}
\begin{list}{$\scriptscriptstyle\bullet$}{
}
\item one or more magnetic or solid-state (SSD) hard disks
 to keep files containing data or programs permanently,
\item (USB) ports for plugging in memory sticks, which make it possible to
keep files
containing data or programs permanently, or to exchange them with other
com\-pu\-ters,
\item a keyboard to enter data or programs into memory,
\item a graphics screen to display texts and images,
\item a mouse to point at places on the graphics screen,
\item a bus to allow fast exchanges of data and
programs between these various components.
\end{list}

When the com\-pu\-ter is switched on, it is wired to load into
its memory and execute a program stored permanently on disk,
called the {\em operating system}\index{Operating system}. The essential
task of this program is to manage all the components of the com\-pu\-ter and to
dialogue with the user by means of a {\em command interpreter}\index{Command interpreter},
which waits for the user to give it orders through
the keyboard and executes them immediately, one after the other,
until the user gives it the order to switch off
the com\-pu\-ter. Among the commands of the interpreter are, for example,
deleting a file stored on disk or launching the
execution of a program written in machine language (generally
called an {\em application program}\index{Application program}) such as, for example, a {\em
text editor} to create a text and store it in a file,
or an {\em interpreter} to translate a \python\ program into machine language and have it
executed.

The first interfaces between man and machine were purely textual. You typed on
the keyboard and saw the result on a printer, soon replaced by a screen with only
lines of text at most 80 characters long. Black-and-white screens were
replaced by colour screens (as can still be seen in the \texttt{terminal} application).
Then came graphics screens made up of very small coloured pixels.
Windowing gives the illusion of several graphics screens on a single one. Then came
the idea of icons to represent files and applications, on which you can
double-click to launch an application. For example, a double click on the icon of
\texttt{LeRobotPythonEN.app} launches the interactive drawing program, whereas dragging the icon
of a program of the \robot\ onto this icon of the application \texttt{LeRobotPythonEN.app} runs this program,
making the drawing window and the terminal appear. Note that you must click on
one or the other of the windows to use it. You can stop the program with Control+C or
ctrl+C (which computer scientists write \^{}C) after selecting the terminal window, or by closing this window, or else with the escape key (\texttt{esc}) in the graphics window.

% (pas d'espace en bas de ce cadre : il tient dans la page)
\end{minipage}
}
%%%%%%%%%%%%%%%%%%%%%%%%%%%%%%%%%%%%%%%%%%%%%%%%%%%%%%%%%%%%%%%%%%%%%%%%%%%

%%%%%%%%%%%%%%%%%%%%%%%%%%%%%%%%%%%%%%%%%%%%%%%%%%%%%%%%%%%%%%%%%%%%%%%%%%%
\noindent \framebox[\textwidth]{
\begin{minipage}{\myparboxwidth}
\setlength{\parindent}{2em}
\small \sf
\vspace{\abovedisplayskip}

Com\-pu\-ters are {\em binary} in the sense that all data, in
particular programs, are represented by sequences of 0s and
1s. The com\-pu\-ter works by transferring an instruction of the program
from a cell of its memory
into its arithmetic and logic unit (via the bus) and
executing it (which in general modifies the content of the memory,
of the disk, of the screen, etc.), then starting again in the same way with the instruction
stored in the next cell of its memory. Some instructions,
called {\em branch instructions}, make it possible to indicate that the next instruction
is stored in the memory cell with a given number.


To execute a program written in machine language and stored on disk,
the operating system asks the arithmetic and
logic unit to load this program into memory, that is, to
copy the program into a part of the memory
that was unoccupied, and then to execute a branch instruction
to the first instruction of this program. The program must end
with a branch instruction back to the operating
system. Modern com\-pu\-ters also have the abil\-i\-ty
to interrupt the execution of the current program and then to resume it
later as if nothing had happened. In the meantime, it is
the operating system that runs, to react according to
the cause of the interruption. For example if the user has moved the mouse,
its new coordinates are stored in memory. If the clock
ticks, which happens every sixtieth of a second for example, the time
is updated, etc. In particular the user has
buttons or keys (%touche Esc (touche d'echappement)
\begin{picture}(14,12)(0, 2)
\put(0,0){\framebox(14,12)[lb]{}}
\put(1.5,2){\makebox(0,0)[lb]{\tiny Esc}}
\end{picture}\hspace{0.3ex},
%touche §
\begin{picture}(12,12)(0, 2)
\put(0,0){\framebox(12,12)[lb]{}}
\put(1.5,2){\makebox(0,0)[lb]{\tiny \S}}
\end{picture}\nolinebreak\hspace{0.3ex}\nolinebreak,
%touche #
\begin{picture}(12,12)(0, 2)
\put(0,0){\framebox(12,12)[lb]{}}
\put(1.5,2){\makebox(0,0)[lb]{\tiny \#}}
\end{picture}\hspace{0.3ex}
 for our \robot) to interrupt
the execution of the program if he or she notices an error.


Among the applications that the operating system can run,
there is one that particularly interests us: the \python\ interpreter.

Modern programming environments (such as {\tt IDLE}, supplied with \python)
take care of editing the text of the
program, that is, of creating and modifying it
by interaction with the user. The
\python\ program is displayed on the screen and stored in memory
as
the user types it on the keyboard. The user can modify it
using various function keys and the mouse. The user is also
given the possibility of saving the text on the hard disk
in a file whose name must be given.

\vspace{\abovedisplayskip}
\end{minipage}
}
%%%%%%%%%%%%%%%%%%%%%%%%%%%%%%%%%%%%%%%%%%%%%%%%%%%%%%%%%%%%%%%%%%%%%%%%%%%

%%%%%%%%%%%%%%%%%%%%%%%%%%%%%%%%%%%%%%%%%%%%%%%%%%%%%%%%%%%%%%%%%%%%%%%%%%%
\noindent \framebox[\textwidth]{
\begin{minipage}{\myparboxwidth}
\setlength{\parindent}{2em}
\small \sf
\vspace{\abovedisplayskip}

When text editing is finished, the environment can
have the program translated into an equivalent program
written in machine language. For so-called {\em compiled}\index{Compilation} languages,
such as \ocaml\ or C, this translation is done once and for all by
a {\em compiler}\index{Compiler}: the compiled program is stored
on disk, then the user goes back to the operating system level to
load into memory
and execute the compiled program as explained above.

During the debugging phase of the program, it is faster to
store the translated program directly in memory, the translator
itself taking care of its execution and then taking back control when
this execution is finished or interrupted. This is how
the \python\ interpreter proceeds, every time it is asked to run a program.

\vfill

One can see the considerable progress that has been made since the
first com\-pu\-ters, for which the wiring had to be modified to
change the program \cite{levy1}:
\begin{list}{$\scriptscriptstyle\bullet$}{
}
\item since the program is stored in memory (following the idea of
Von Neumann\footnote{\ \href{https://en.wikipedia.org/wiki/John_von_Neumann}{John von Neumann}, 1903--1957, American mathematician.})
and the memory is connected to fast disks, it is possible to
change the program of the machine almost \hbox{instantly;}
\item the use of high-level languages such as \python\ or \ocaml\ and of efficient interpreters
or compilers
makes it possible to avoid having to program in a thankless machine
\hbox{language;}
\item the alphanumeric screen lets the user edit his or her program
interactively at the keyboard, which saves a considerable amount of time
compared with handling the cardboard punched cards
still used in the \hbox{1970s;}
\item computing power, linked to the density of transistors on
computer chips (which ranges from a few billion to several hundred billion
transistors in 2016), makes it possible to manage graphics
screens directly connected to the memory, which allows
very fast changes and hence the use of windows, mice, \hbox{etc.;}
\item the size of the memory (a few billion characters, GB, gigabytes) and
of the disks (a few thousand billion characters, TB, terabytes) makes it possible to consider
ever larger programs and data, and therefore ever more
sophisticated applications.
\end{list}
\vspace{\abovedisplayskip}
\end{minipage}
}
%%%%%%%%%%%%%%%%%%%%%%%%%%%%%%%%%%%%%%%%%%%%%%%%%%%%%%%%%%%%%%%%%%%%%%%%%%%

%%%%%%%%%%%%%%%%%%%%%%%%%%%%%%%%%%%%%%%%%%%%%%%%%%%%%%%%%%%%%%%%%%%%%%%%%%%
\noindent \framebox[\textwidth]{
\begin{minipage}{\myparboxwidth}
\setlength{\parindent}{2em}
\small \sf
\vspace{\abovedisplayskip}
The size of the largest programs that can be written today
now runs to tens or hundreds of millions of lines, which is considerable
and not without difficulties, since the limits of human
understanding are being reached. A training of the mind, a
new rigour, even more precise than the logic of mathematicians,
becomes the prime quality which, together with imagination and a sense of
aesthetics, makes good computer scientists.

\indent A little method also
helps a lot. The whole history of programming languages
is full of inventions making it possible to formulate more clearly how
to organize data in memory or the computations of the com\-pu\-ter.
Each time, human capacities, and therefore the maximum size of the largest
 programs, have been pushed back.
A very old idea for structuring computations, due to Von Neumann and Turing%
\footnote{\ \href{https://en.wikipedia.org/wiki/Alan_Turing}{Alan Turing}, 1912--1954, English mathematician.}, is that
of the {\em procedure}, formerly called a {\em routine}
or {\em subprogram}. We begin to introduce it in the next chapter.
\vspace{\abovedisplayskip}
\end{minipage}
}
%%%%%%%%%%%%%%%%%%%%%%%%%%%%%%%%%%%%%%%%%%%%%%%%%%%%%%%%%%%%%%%%%%%%%%%%%%%
\newpage
\thispagestyle{empty}
